\section{Literature review}
\label{sec:literature}
\subsection{Predictive information and the timing of execution}
Optimal execution studies how to complete an order while balancing the consequences of trading quickly against those of carrying inventory. The impact--risk formulation of \citet{almgren2001} supplies a foundational benchmark. Once prices contain predictive information, the schedule also depends on when favourable prices are expected to occur. \citet{cartea2016} incorporate expected future net order flow, while \citet{bechler2015} study execution with dynamic order-flow imbalance. \citet{cartea2018} connect order-book signals to trading decisions. These studies motivate using an imbalance predictor, but do not establish that the coarser taker-volume measure available in minute bars contains the same information as an observed order book.

The relevant forecasting object is generally a temporal profile. \citet{lehalle2019} incorporate signals into execution and derive deterministic schedules conditional on initial information within part of their analysis. \citet{neuman2022} instead address signal-adaptive trading with temporary and transient impact. Learning latent alpha states provides another route for incorporating information into trading controls \citep{casgrain2019}. In the related portfolio-choice setting, \citet{garleanu2013} show why transaction costs and the persistence of return predictors affect dynamic adjustment. These formulations differ in information structure and feasible actions, so they should not be treated as interchangeable benchmarks for a schedule fixed at the start of an episode.

Two contributions are particularly close to the distinction developed here. \citet{forde2022} express predictive execution through expected remaining price changes in a Gaussian-signal model with power-law resilience. \citet{abijaber2025} use conditional price contrasts in a more general propagator framework. Their formulations already contain the economic reason that fixed-volume execution depends on relative prices across trading times. Recent work by \citet{bank2026} further develops execution and speculation with trade signals. Consequently, neither using forecasts in execution nor cancelling a common price-level component constitutes a new execution principle. Our narrower contribution is an observable diagnostic: holding the terminal forecast exactly fixed while changing the intrahorizon profile separates information captured by terminal prediction scores from information that can alter a schedule.

\subsection{Decision loss, calibration and model uncertainty}
The distinction between prediction quality and decision quality extends well beyond execution. \citet{elmachtoub2022} formulate learning around the loss induced by downstream optimization. \citet{liu2021} establish risk and calibration results for a related predict-then-optimize method under stated assumptions. The review and benchmark of \citet{mandi2024} place these approaches within decision-focused learning, while \citet{gupta2024} develop directional-gradient methods. Dependence in the training observations creates additional theoretical issues, as the preprint of \citet{liudependent2024} illustrates. Our analysis shares the motivation of evaluating forecasts through decisions, but does not train a new decision-focused predictor or establish comparable generalization bounds. Ordinary least squares is retained so that the decision diagnostic can be examined without attributing results to a complex forecasting architecture.

Economic forecast evaluation also appears in applied financial research. \citet{maciel2025}, for example, evaluates equity risk-premium forecasts through a portfolio allocation strategy. That application concerns investment exposure rather than the timing of a committed order. The distinction matters: an investor can change how much risk to hold, whereas every feasible schedule here must complete the same sale. A forecast of the final price may inform exposure choice while failing to distinguish competing execution profiles. The matched-endpoint comparison makes this restriction explicit, and the quadratic decomposition identifies both the gross price-alignment benefit and the extra scheduling cost.

Attenuating uncertain signals is likewise an established concern. \citet{kan2007} study portfolio choice under parameter uncertainty. \citet{cartea2017} address model uncertainty in algorithmic trading, and the liquidation working paper of \citet{cartea_liquidation2017} focuses on ambiguity aversion. \citet{horst2022} examine liquidation under factor uncertainty, while \citet{demarch2018} study signal-based trading with uncertainty and threshold behaviour. These approaches motivate caution about an estimated signal, but their objectives and uncertainty sets differ from a coefficient fitted on a preceding calibration month. We therefore interpret our coefficient as an empirical schedule weight, not as an ambiguity-robust optimum.

Uncertainty about execution costs is distinct from uncertainty about returns. \citet{fouque2022} jointly consider signals and stochastic price impact. The preprint of \citet{hey2023} examines the cost of misspecifying impact, and \citet{neumanoffline2023} develop an offline-learning approach to propagator models. \citet{macri2024} study reinforcement learning when liquidity varies over time. This strand limits what can be inferred from the present exercise: changing an assumed impact parameter is a sensitivity analysis, not an estimate of the impact that a real order would generate. A transparent small model helps locate the source of a replay result, but cannot substitute for measurements of depth, fills or endogenous price responses.

\subsection{Uncertainty gates and dependence-aware evaluation}
A particularly relevant recent comparator is \citet{irshad2026}, whose uncertainty-aware execution framework uses a conformal-prediction gate and a baseline fallback. The idea of declining to act on an uncertain prediction therefore precedes this study. Baseline protection also appears in reinforcement learning through safe policy improvement \citep{laroche2019}. Our fixed convex mixture remains feasible and can equal the baseline, but this alone provides no statistical guarantee that its future cost is lower. The empirical result that the conservative calibration selects zero in every window must be read as abstention, rather than evidence that a new trading rule improves execution.

Prediction-interval coverage and economic protection are separate properties. \citet{gibbs2021} develop adaptive conformal inference under distribution shift, and \citet{barber2023} analyse conformal prediction beyond exchangeability. \citet{zaffran2022} and \citet{xu2021} address time-series settings through different constructions and assumptions. These results explain why an uncertainty procedure for dependent data needs an explicit target and justification; they do not confer distribution-free protection on a bootstrap schedule-weight rule. No conformal guarantee is claimed here. The relevant empirical questions are whether historical price alignment compensates for the stated costs and whether that relation persists in the next evaluation month.

Evaluation also requires separating temporal dependence from research selection. \citet{politis1994} provide the stationary bootstrap as one method for dependent observations. Our calendar-day cluster bootstrap is a different procedure: it preserves within-day episode grouping and simultaneous observations for the two assets, but does not preserve dependence across days or repeat model estimation in each resample. \citet{bailey2017} analyse backtest overfitting, and \citet{arnott2019} advocate disciplined backtesting protocols. These concerns motivate reporting every specified asset-month and cost scenario, retaining the original analysis chronology, and identifying the expanded study as exploratory. Chronological fitting prevents future observations from entering an episode's estimated rule; it does not undo knowledge obtained when designing the wider study.

\subsection{Cryptocurrency context and the research gap}
Cryptocurrency is already an established execution application. \citet{bundi2024} study optimal trade execution in cryptocurrency markets, while \citet{schnaubelt2022} examine limit-order placement using deep reinforcement learning. The latter problem depends on execution opportunities and order-book information that cannot be reconstructed from minute bars. \citet{barucci2023} investigate impact and efficiency in cryptoasset markets; \citet{almeida2024} review the wider microstructure literature. Exchange segmentation documented by \citet{makarov2020} further cautions against generalizing a single-exchange result to the asset class. These studies support an exchange-specific, resolution-specific interpretation of the Bitcoin and Ether replay.

The resulting research question is deliberately bounded. Can terminal-equivalent forecasts produce different feasible schedules and economic outcomes, and can a simple benefit--cost decomposition explain the difference? Existing theory supplies the underlying price contrasts, and existing decision-focused research supplies the distinction between forecasting and optimization losses. The contribution here is to connect those ideas in a reproducible matched-endpoint experiment, report their implications across rolling cryptocurrency samples, and identify when favourable average price alignment is outweighed by the assumed scheduling costs. It is a diagnostic contribution rather than evidence of a newly profitable forecasting or execution strategy.
